\section{Билет 6. Математический маятник. Уравнение малых колебаний и период колебаний математического маятника с выводом.}

\begin{definition}
\textbf{Математическим маятником} называется идеализированная механическая система, состоящая из материальной точки массы $m$, подвешенной на невесомой и нерастяжимой нити длины $l$ в однородном поле тяжести. На практике хорошим приближением служит небольшой тяжёлый шарик на длинной тонкой нити.
\end{definition}

\begin{theorem}[Уравнение движения математического маятника]
Пусть маятник отклонён от вертикали на угол $\varphi$. Сила тяжести $m\vec{g}$ создаёт вращательный момент относительно точки подвеса, стремящийся вернуть маятник в положение равновесия:
\begin{equation}
    N = -mgl \sin\varphi, \label{eq:torque_06}
\end{equation}
где знак «минус» указывает на то, что момент направлен в сторону, противоположную углу отклонения.

Запишем основное уравнение динамики вращательного движения:
\begin{equation}
    I \ddot{\varphi} = N,
\end{equation}
где момент инерции материальной точки относительно точки подвеса $I = ml^2$. Подставляя \eqref{eq:torque_06}, получаем:
\begin{equation}
    ml^2 \ddot{\varphi} = -mgl \sin\varphi \quad \Rightarrow \quad \ddot{\varphi} + \frac{g}{l} \sin\varphi = 0. \label{eq:exact_pendulum}
\end{equation}
Это нелинейное дифференциальное уравнение. Для малых углов отклонения ($\varphi \lesssim 0.1$ рад $\approx 5.7^\circ$) справедливо приближение $\sin\varphi \approx \varphi$. Тогда уравнение принимает вид линейного однородного ДУ второго порядка:
\begin{equation}
    \ddot{\varphi} + \frac{g}{l} \varphi = 0. \label{eq:small_pendulum}
\end{equation}
Обозначая $\omega_0^2 = \frac{g}{l}$, приходим к каноническому уравнению гармонических колебаний:
\begin{equation}
    \ddot{\varphi} + \omega_0^2 \varphi = 0.
\end{equation}
\end{theorem}

\begin{property}
Решением уравнения \eqref{eq:small_pendulum} является закон движения $\varphi(t) = \varphi_{\max} \cos(\omega_0 t + \varphi_0)$. Период собственных малых колебаний математического маятника равен:
\begin{equation}
    T = \frac{2\pi}{\omega_0} = 2\pi \sqrt{\frac{l}{g}}. \label{eq:period_pendulum}
\end{equation}
Из формулы \eqref{eq:period_pendulum} следует, что период:
\begin{itemize}
    \item \textbf{Не зависит от массы} груза $m$.
    \item \textbf{Не зависит от амплитуды} колебаний (при условии малости углов --- изохронность).
    \item Определяется исключительно длиной нити $l$ и ускорением свободного падения $g$ в данной местности.
\end{itemize}
\end{property}

\begin{remark}
При больших амплитудах приближение $\sin\varphi \approx \varphi$ нарушается, и колебания перестают быть строго гармоническими. Период начинает зависеть от амплитуды $\varphi_{\max}$ и вычисляется через эллиптические интегралы. Разложение в ряд даёт первую поправку: $T \approx 2\pi\sqrt{l/g} \left(1 + \frac{1}{16}\varphi_{\max}^2 + \dots\right)$. В рамках курса общей физики всегда используется формула \eqref{eq:period_pendulum}.
\end{remark}

\begin{figure}[htbp]
    \centering
    \begin{tikzpicture}[
        scale=1.5,
        pivot/.style={fill=gray!50, draw=black, thick},
        rod/.style={thick, black},
        mass/.style={fill=blue!30, draw=blue!70, thick},
        force/.style={->, >=Stealth, thick},
        angle/.style={->, >=Stealth, blue!70!black}
    ]
    % Точка подвеса
    \fill[gray!30] (-1.2, 3) rectangle (1.2, 3.2);
    \draw[thick] (-1.2, 3) -- (1.2, 3);
    \draw[pivot] (0, 3) circle (3pt);
    \coordinate (O) at (0, 3);

    % Вертикальная линия (пунктир)
    \draw[dashed, gray] (0, 3) -- (0, -0.5);
    \coordinate (V) at (0, 0); % Положение равновесия

    % Маятник в отклоненном положении
    \def\angle{30} % Угол отклонения
    \def\length{3} % Длина нити

    % Координата груза
    \coordinate (M) at ({\length*sin(\angle)}, {3 - \length*cos(\angle)});

    % Нить
    \draw[rod] (O) -- (M);

    % Груз
    \draw[mass] (M) circle (0.25);
    \node at (M) {$m$};

    % Длина нити l
    \draw[<->, gray] (-0.3, 3) -- (-0.3, 0) node[midway, fill=white, font=\small] {$l$};

    % Угол отклонения theta
    \draw[angle] (0, 2.3) arc (90:90-\angle:0.7);
    \node at (0.4, 2.5) {$\theta$};

    % Силы
    % Сила тяжести mg (вниз)
    \draw[force, red] (M) -- ++(0, -1.2) node[midway, right, red] {$mg$};

    % Сила натяжения нити T (вверх вдоль нити)
    \draw[force, blue!70!black] (M) -- ++({-0.4*sin(\angle)}, {0.4*cos(\angle)})
        node[midway, left, blue!70!black] {$T$};

    % Касательная составляющая силы (возвращающая сила)
    \draw[force, green!60!black] (M) -- ++({-1.0*cos(\angle)}, {-1.0*sin(\angle)})
        node[midway, below left, green!60!black] {$F_{\tau}$};

    % Траектория движения (дуга)
    \draw[dashed, gray!50] ({-\length*sin(20)}, {3 - \length*cos(20)})
        arc (250:290:\length);

    % Ось координат
    \draw[->, thick] (-2, -0.5) -- (2, -0.5) node[right] {$x$};
    \node at (0, -0.7) {$O$};

    % Подписи
    \node[anchor=west, align=left, font=\small] at (1.5, 2.5)
        {\textbf{Параметры:}\\[3pt]
        $l$ --- длина нити\\[2pt]
        $m$ --- масса груза\\[2pt]
        $\theta$ --- угол отклонения};

    \end{tikzpicture}
    \caption{Математический маятник: груз массы $m$ на невесомой нерастяжимой нити длины $l$. Возвращающая сила $F_{\tau} = -mg\sin\theta$ направлена по касательной к траектории.}
    \label{fig:mathematical_pendulum}
\end{figure}
